\section{Introduction}\label{intro:overview}

The Penrose inequality compares the mass seen at infinity with the size
of a region hidden behind a trapped boundary. An initial data set
consists of a Riemannian manifold $(M,g)$ and a symmetric two-tensor $K$,
representing the metric and second fundamental form of a spacelike
hypersurface. For a hypersurface $S\subset M$ and a unit normal $\nu$
toward the chosen exterior, its future outward null expansion is
\[
                 \theta_+(S)=H_S+\tr_S K,
\]
where $H_S$ is positive on Euclidean spheres with outward normal.
We call $S$ weakly future outer trapped when $\theta_+(S)\le0$.
A future marginally outer trapped surface (MOTS) satisfies $\theta_+=0$.
The dominant energy condition requires the energy density of the
constraint equations to be at least the norm of their momentum density.
Precise normalizations and asymptotic assumptions are given separately
for each spatial dimension below. The asymptotically flat results resolve
the minimum-enclosing-area form of the spacetime Penrose conjecture in
three and four spatial dimensions under those hypotheses.

The geometric quantity throughout the asymptotically flat results is
an enclosing infimum. An enclosing cut separates the entire designated
inner boundary from the chosen distant end. Its full area is counted,
including portions coincident with the original boundary. Cuts may be
disconnected and need not themselves be trapped. This distinction matters:
a trapped surface need not minimize area among enclosing surfaces.
When several ends are present, the definition also specifies which of
them lie on the excluded side. We give the exact admissible classes in
each theorem, rather than identifying infima with different end constraints.
In this overview, $A_{\min}$ denotes the enclosing infimum appropriate
to the result being described; each part defines its precise cut class.

\subsection{The results and their relation}

For an asymptotically flat end, let $(E,P)$ denote its ADM energy and
momentum. Once $E>|P|$, its invariant mass is
$m=\sqrt{E^2-|P|^2}$. In three spatial dimensions, the basic numerical
bound is
\[
                     m\ge\sqrt{\frac{A_{\min}}{16\pi}}.
\]
Here $A_{\min}$ denotes the relevant minimum enclosing area.
The one-ended theorem is proved first under a timelike assumption on
$(E,P)$. A positive conformal deformation then rules out a non-timelike
ADM vector, and cutting off the unwanted distant ends gives the
complete-data formulation. The latter permits a disconnected chosen
exterior region as well as a disconnected trapped boundary.
The precise statements are
Theorem~\ref{n3:intro:exterior-inequality},
Corollary~\ref{bridge:bridge:timelike}, and
Theorem~\ref{bridge:intro:global-penrose}.

Equality concerns the original data and requires additional geometric
hypotheses. In Theorem~\ref{n3:intro:rigidity}, the chosen exterior is
connected and one-ended, and its finite marginally trapped boundary
is outer area-minimizing. Its outermostness condition excludes an
enclosing weakly trapped cut that replaces any nonempty subset of the
boundary components, even if other components remain coincident.
Equality then forces a single spherical boundary component and a
Schwarzschild spacetime realization of the entire exterior. The proof
does not prescribe the data on the other side of that boundary.

For electric and magnetic fields, write $Q_E,Q_B$ for their total
charges, $Q=(Q_E^2+Q_B^2)^{1/2}$, and
$r=(A_{\min}/4\pi)^{1/2}$. Theorem~\ref{ch:intro:main} proves
\[
             m\ge Q,\qquad r\le m+\sqrt{m^2-Q^2}
\]
on a one-ended exterior. Its matter condition includes the
electromagnetic momentum density. It allows arbitrary $K$ and nonzero
ADM momentum. The familiar lower bound
$m\ge(r+Q^2/r)/2$ follows on the branch $r\ge Q$; it is not an equivalent
formulation for every positive $r$. Purely electric rest-frame
consequences, including a designated-end theorem with finitely many
ends, follow with their own explicit cut convention. No charged
equality classification is asserted.

In four spatial dimensions, an enclosing hypersurface has
three-dimensional volume $A_e$. With $\omega_3=2\pi^2$, the theorem is
\[
              m_e\ge\frac12\left(\frac{A_e}{\omega_3}\right)^{2/3}.
\]
Theorem~\ref{four:thm:main} allows arbitrary second fundamental form,
finitely many asymptotically flat ends, and disconnected weakly future
outer trapped boundary. Its hypotheses impose neither symmetry nor
maximality. Theorem~\ref{rig4:thm:rigidity} completes this bound in the
connected, one-ended horizon class: if the boundary is a future MOTS,
outermost and outer area-minimizing, equality holds exactly for original
exterior data realized by a smooth spacelike hypersurface in the regular
Schwarzschild--Tangherlini extension. The embedding recovers both $g$ and
$K$, reaches spatial infinity, and attaches smoothly to a future-horizon
section or the bifurcation sphere. Boosted ends and non-time-symmetric
slices are included. These additional horizon assumptions belong only
to the equality theorem; no disconnected equality or behind-horizon
classification is asserted in four dimensions.

The last part treats a different asymptotic geometry. For cosmological
constant $-3$, fix a positive-mass Schwarzschild--anti-de Sitter
background, a transverse-traceless tensor, and a decaying positive
conformal solution branch generated by that tensor. Here
``transverse-traceless'' means divergence free and trace free in the
background metric. Theorem~\ref{ads:thm:main} gives, on a sufficiently
small interval depending on these fixed choices,
\[
 m_{\AH}\ge\sqrt{\frac{A}{16\pi}}
                       \left(1+\frac{A}{4\pi}\right).
\]
The quantity $m_{\AH}$ is the positive Lorentz norm of the hyperbolic
mass covector, and $A$ is the area of the prescribed marginally trapped
boundary. The theorem establishes equality precisely for radial seeds
at sufficiently small positive parameters. Identifying $A$ with an
enclosing infimum is a separate corollary with additional geometric
hypotheses. The interval is not uniform over seeds or branches.

\subsection{How the proofs fit together}

The asymptotically flat numerical proofs have a common structure.
First, the distant end is replaced and repaired so that its energy
approaches the original invariant mass in a rest frame. The replacement
preserves the dominant energy condition and controls every enclosing
cut. Next, trapped-region boundaries and disjoint collars prepare a
fixed exterior for a coupled elliptic deformation. The unknown graph
function and conformal factor produce a Riemannian comparison metric
with nonnegative scalar curvature, a controlled energy change, and an
enclosing area no smaller than the required original infimum.
The Riemannian Penrose inequality completes the comparison.

The three-dimensional elliptic construction is proved in full once.
The charged argument uses this same system: it verifies its hypotheses,
transports closed flux forms through the geometric changes, and proves
a square estimate incorporating the electromagnetic momentum.
In four dimensions the scalar identity and regularity estimates
change, so we give the corresponding construction separately.
Its final metric has zero second fundamental form. A minimizing
enclosure therefore brings that actual metric directly within the
Riemannian theorem's hypotheses.

The equality proofs start from the unmodified exterior. Variations of
the constraints and of the enclosing-area infimum give a normalized
causal stationary field. A twist estimate makes its positive-norm
region static, without assuming regularity of an interior zero set.
A complete conformal double then forces Euclidean geometry. In three
dimensions, the zero-mass argument uses the spin structure supplied
by orientability, and the boundary identities give each component
the total original area, forcing one component. In four dimensions,
we prove the required compactification and reduce zero-mass rigidity
to the smooth positive-mass theorem with arbitrary other ends.
The final step in both cases reconstructs the original $g,K$ and
the smooth horizon attachment; it does not transfer equality through
the numerical deformations. Figure~\ref{intro:proof-routes} shows
these two uses of the inequality.

Two ingredients keep the original-data argument independent of a
chosen minimizing surface. Full obstacle perimeter includes every
boundary-coincident piece. Compactness of all minimizing enclosures
then gives the one-sided metric derivative as the minimum of their
area derivatives. A positive measure on this minimizing space yields
the variational identity; normal tests and the precise outermostness
condition concentrate its area term on the original horizon.
The causal-field estimates also separate a reusable finite-jet
asymptotic argument from the charge normalization, so the converse
can establish its asymptotic spacelike plane without assuming it.

\begin{figure}[tbp]
\centering
\begin{tikzpicture}[
  box/.style={draw=black!45,rounded corners=2pt,align=center,
    text width=6.2cm,minimum height=14mm,inner sep=5pt,font=\small},
  arr/.style={-{Stealth[length=2mm]},thick}]
\node[box] (data) at (0,0) {Original data with DEC\\and trapped boundary};
\node[box] (deform) at (0,-2.2) {End reduction and deformation\\to a Riemannian comparison metric};
\node[box] (bound) at (0,-4.4) {Riemannian Penrose inequality\\and passage to the original bound};
\node[box] (equal) at (7.2,0) {Equality on the original data\\with the stated horizon hypotheses};
\node[box] (field) at (7.2,-2.2) {Area and constraint variations\\give a causal stationary field};
\node[box] (recover) at (7.2,-4.4) {Static classification and recovery\\of the original spacelike hypersurface};
\draw[arr] (data)--(deform);
\draw[arr] (deform)--(bound);
\draw[arr] (equal)--(field);
\draw[arr] (field)--(recover);
\end{tikzpicture}
\caption{The numerical comparison and the equality argument. The bound
proved on the left applies to the nearby data used in the variations on
the right. Rigidity is obtained from the original data.}
\label{intro:proof-routes}
\end{figure}


The anti-de Sitter proof uses the Taylor coefficients of the given
branch. Actual weighted remainder bounds connect those coefficients
to the mass and area. The quadratic coefficient is a nonnegative quadratic
form with a kernel of dimension four. A fourth-order argument treats its
nonradial directions,
while radial directions satisfy an exact mass conservation identity.
This part is independent of the asymptotically flat elliptic construction.

\subsection{Relation to earlier work}

Penrose proposed the mass--horizon comparison in the context of
gravitational collapse and cosmic censorship \cite{Penrose1973}.
For time-symmetric data, Huisken and Ilmanen proved the
connected-horizon inequality by weak inverse mean curvature flow,
and Bray proved the total-area bound for possibly disconnected
horizons by conformal flow \cite{HuiskenIlmanen2001,Bray2001}.
Bray and Lee extended the numerical Riemannian inequality to spatial
dimensions below eight; spin enters their higher-dimensional equality
statement \cite{BrayLee2009}. Ben-Dov's counterexample distinguishes
apparent-horizon area from minimum enclosing area \cite{BenDov2004}.
Carrasco and Mars constructed counterexamples for outermost generalized
apparent horizons \cite{CarrascoMars2010}.
Bray and Khuri developed a generalized Jang reduction, including
spherical existence results and conditional nonspherical constructions
\cite{BrayKhuri2010,BrayKhuri2011}.

Allen, Bryden, Kazaras, and Khuri obtained a universal suboptimal
mass--area bound in three dimensions under their decay and horizon
assumptions \cite{AllenBrydenKazarasKhuri2025}.
Khuri and Kunduri proved the sharp inequality for
$SU(n+1)$-invariant cohomogeneity-one data in spatial dimensions
$2(n+1)$ \cite{KhuriKunduri2025}.
The inequalities here have a different scope: the three- and
four-dimensional results impose no such symmetry, and use their
specified minimum enclosing quantities.

The stationary field arising from equality is related to the
variational study of ADM mass. Hirsch and Huang proved a strong
maximum principle for the norm of a causal Killing field and studied
mass minimizers with fixed Bartnik boundary data
\cite{HirschHuang2025}. Those boundary data and regularity hypotheses
differ from the area variations used here. Once static vacuum geometry
is obtained, the conformal-doubling method of Bunting and
Masood-ul-Alam provides the classical model for the classification
\cite{BuntingMasoodUlAlam1987}. We prove the completeness and
low-regularity zero-mass statement needed for our double.
Higher-dimensional static-vacuum uniqueness was developed by Gibbons,
Ida, and Shiromizu \cite{GibbonsIdaShiromizu2002}. Here the static geometry
and its complete base must first be recovered from the original initial data.
The charged and anti-de Sitter parts give the comparisons with
earlier results specific to those settings where their hypotheses
can be stated precisely.

\subsection{Reading conventions}

The five parts are ordered by their mathematical dependencies.
The three-dimensional numerical theorem precedes the global extension,
rigidity, and charged application. The four-dimensional part states its own normalizations and contains
both the numerical construction and original-data equality proof.
The anti-de Sitter part is independent and uses its own asymptotic
geometry. The entry points are
Sections~\ref{n3:intro:section}, \ref{n3:rigidity:setup},
\ref{ch:intro:charged-section}, \ref{four:sec:introduction}, and
\ref{ads:sec:setting}, respectively.
Within each analytic construction the geometry
and strictification are fixed before the deformation parameters.
Uniform polynomial estimates and regularity estimates at a fixed
parameter have distinct roles; the proofs specify the order of
exhaustion and limits where each is used.
